\begin{helper}[FiberAndDegreeMixedLiftedIntersectionUniformBound]
\label{helper:FiberAndDegreeMixedLiftedIntersectionUniformBound}
Let \(m=\noderef{TwoBiteNaturalReducedVertexCount}(n)\) and
\[
p=\noderef{TwoBiteEdgeProbability}(1/2,n).
\]
Assume \(n\le m^2\) and \(\log n\ge1\).  For fixed base graphs \(G_R\) and
\(G_B\) on \(\operatorname{Fin}(m)\), suppose that for every
\(r,b\in\operatorname{Fin}(m)\),
\[
|\{u:G_R.Adj(r,u)\}|\le2pm
\quad\text{and}\quad
|\{v:G_B.Adj(b,v)\}|\le2pm .
\]
Then the \(\noderef{TwoBiteConditionalProbability}\), conditioned on
\(\omega.1=G_R\) and \(\omega.2.1=G_B\), of the event that some mixed lifted
intersection has size greater than \(100(\log n)^3\) is at most
\[
m^2\exp(-2(\log n)^3).
\]
\end{helper}

\begin{proof}
Write \(L=\log n\), \(N=m^2\), and \(T=100L^3\).  Since \(L\ge1\), we have
\(T>0\).  We first identify the conditional distribution after fixing the
two base graphs.  By \(\noderef{TwoBiteSample}\), a sample is a triple
\[
\omega=(G_R',G_B',\phi),
\]
where \(\phi\) is an injection from the \(n\) final vertices into
\(\operatorname{Fin}(m)\times\operatorname{Fin}(m)\).  The conditioning event
in \(\noderef{TwoBiteConditionalProbability}\) is exactly
\[
G_R'=G_R,\qquad G_B'=G_B.
\]
By \(\noderef{TwoBiteSampleWeight}\), the weight of such a triple factors as
the red graph weight, the blue graph weight, and
\(\noderef{UniformInjectionWeight}(n,m)\).  Inside this conditioning cell the
two graph weights are constant and cancel between numerator and denominator.
The assumption \(n\le m^2\) makes the injection space nonempty, so
\(\noderef{UniformInjectionWeight}\) is the reciprocal of the number of
injections.  Thus, conditional on \(G_R\) and \(G_B\), the injection is
uniform over all injections.  Equivalently, its image is a uniform
\(n\)-element subset sampled without replacement from the \(N=m^2\) product
cells: every \(n\)-set of cells has exactly \(n!\) orderings by the final
vertices, so all image sets have the same total conditional weight.

Fix \(r,b\in\operatorname{Fin}(m)\).  By
\(\noderef{TwoBiteLiftedNeighborhood}\), a final vertex lies in
\[
\noderef{TwoBiteLiftedNeighborhood}(\omega,\operatorname{inl}(r))
\cap
\noderef{TwoBiteLiftedNeighborhood}(\omega,\operatorname{inr}(b))
\]
exactly when its injected cell lies in the marked rectangle
\[
K_{r,b}=N_{G_R}(r)\times N_{G_B}(b).
\]
The degree hypothesis for this same pair \((r,b)\) gives
\[
|K_{r,b}|
=|N_{G_R}(r)|\,|N_{G_B}(b)|
\le (2pm)(2pm)=4p^2m^2.
\tag{1}
\]
Choose any bijection between the \(m^2\) product cells and
\(\operatorname{Fin}(N)\).  Under this bijection the count of sampled cells
in \(K_{r,b}\) is a hypergeometric count with population \(N=m^2\), sample
size \(n\), and marked set size \(|K_{r,b}|\).  Its mean satisfies
\[
\mu_{r,b}=n\frac{|K_{r,b}|}{m^2}\le 4p^2n.
\tag{2}
\]

Now use the specialized parameter value.  Since
\(p=\noderef{TwoBiteEdgeProbability}(1/2,n)\), the definition of
\(\noderef{TwoBiteEdgeProbability}\) gives
\[
p=\frac12\sqrt{\frac Ln}.
\]
Because \(L\ge1\), \(n>0\) and \(L/n\ge0\), so squaring this identity is
legitimate and yields
\[
4p^2n=L.
\tag{3}
\]
Combining (2) and (3), \(\mu_{r,b}\le L\).

Instantiate \(\noderef{FiberAndDegreeMixedLiftedIntersectionHypergeometricBound}\)
with population \(N=m^2\), sample size \(n\), marked set \(K_{r,b}\),
threshold \(T=100L^3\), and exponential parameter \(L_{\mathrm{param}}=1\).
It gives
\[
\Pr(Y_{r,b}\ge T\mid G_R,G_B)
\le \exp(\mu_{r,b}(e-1)-T).
\tag{4}
\]
Using \(\mu_{r,b}\le L\), \(e-1<2\), and \(L\ge1\),
\[
\mu_{r,b}(e-1)-T
\le 2L-100L^3
\le -2L^3.
\tag{5}
\]
Thus, for this fixed pair,
\[
\Pr(Y_{r,b}>100L^3\mid G_R,G_B)
\le \exp(-2L^3),
\tag{6}
\]
because the event \(Y_{r,b}>100L^3\) is contained in
\(Y_{r,b}\ge T\).

There are exactly \(m^2\) choices of the mixed pair \((r,b)\).  Applying the
union bound to (6) over all such choices gives
\[
\Pr(\exists r,b:\,Y_{r,b}>100L^3\mid G_R,G_B)
\le m^2\exp(-2L^3),
\]
which is the asserted conditional bound.
\end{proof}
